\subsection{Limitaciones y evaluación del aprendizaje obtenido}
\label{sec:limites-rl}
La recompensa representa una aproximación al objetivo que se busca mejorar, por lo que maximizarla puede favorecer comportamientos que aprovechen errores del evaluador sin aumentar la calidad de las respuestas. Este fenómeno suele denominarse reward hacking, mientras que la sobreoptimización describe el deterioro que puede aparecer al optimizar excesivamente una señal imperfecta. Gao y sus colaboradores estudiaron esta divergencia mediante un entorno sintético con un modelo de recompensa de referencia y otro utilizado como aproximación, mostrando que el aumento de la señal optimizada puede coexistir con un deterioro de la evaluación de referencia~\cite{gao2022overoptimization}.

Para evaluar el aprendizaje obtenido se distingue la recompensa de entrenamiento del desempeño sobre consultas reservadas y de los resultados de un benchmark externo. La primera describe el ajuste a la señal utilizada, mientras que las otras permiten examinar si ese ajuste mejora las respuestas fuera del entrenamiento. Esta comparación requiere datos separados y protocolos consistentes, y debe complementarse con la revisión de errores que la recompensa puede pasar por alto. En el caso clínico, recuperar los conceptos esperados no basta para comprobar la pertinencia ni la seguridad de todas las recomendaciones, por lo que las rúbricas y las respuestas también requieren revisión experta.